\begin{abstract}
\noindent
Conversational AI chatbots like ChatGPT and Claude have taken over the kinds
of queries that used to drive search engine advertising. These include areas like product research,
comparison shopping, recommendation requests. Yet, they have no working
advertising format of their own. Banner ads need screen space the chat
surface does not have. Pop ups and interstitials break the back and forth
flow that makes chatbots useful in the first place. Search style keyword
auctions look at only the latest message and miss the context built up over
a multi-turn conversation. The result is a large and growing commercial
surface with no native way to monetise it.
\newline \newline
This report presents \emph{Ad LLM as a Service}. This is a working platform that
lets a large language model weave sponsored product recommendations directly
into its own replies, in a way that feels like a genuine suggestion rather
than an inserted ad. The platform answers the X-as-a-Service brief by
turning this capability of native conversational advertising into a
cloud-served product that any advertiser can plug into. It realises the
dual-utility crowdsourcing principle by design whereby every ad the system serves
does two jobs at once, delivering a commercial impression to the advertiser
and producing a structured record of how the user reacted, which the
platform then uses to automatically improve how future ads are presented.
\newline \newline
The platform is built around eleven components organised into four layers.
The chat pipeline is the core path every user message goes through. A
lightweight model first reads the message to work out what the user wants
and whether an ad would even be appropriate, a session store keeps the last
ten turns of context in memory, a semantic search looks up the most
relevant ads from the inventory, a ranker picks a winner by blending how
well the ad matches, how much the advertiser is willing to pay, and how
close the user seems to be to a purchase, and a second language model
writes the final reply with the chosen ad woven in and a clear
\texttt{[Sponsored]} label attached. Alongside the chat path, an
image-generation surface does the same thing for pictures: when a user
asks for an image, the platform quietly rewrites the prompt to place a
matched product naturally into the scene, and, when the advertiser has
supplied a real photograph of the product, feeds that photograph to the
image model so the rendered version actually looks like the product rather
than a guess at it.
\newline \newline
Behind the scenes, a feedback loop watches what users do next. After every
ad, a judge model reads the follow up message and scores the interaction
along six dimensions. They include things like whether the user engaged at all, how
natural the ad felt, and how close they came to wanting to buy. Those scores 
feed an agentic optimiser that decides on its own when an ad needs attention and 
rewrites its presentation guide when the numbers start slipping. Every 
rewrite is saved with a plain English
explanation of why it was made, so the history of how each ad has been
presented over time is fully auditable. On the commercial side, a
self-serve portal lets advertisers sign up, create campaigns with total and
daily budget caps, and draft ads by simply pasting a product URL that the
system scrapes and fills in for them. Two dashboards, one for the platform
operator, one for each advertiser, report performance live, and because
both read directly from the same event log, the numbers on the screen
always match the data driving the feedback loop.
\newline \newline
The user's answer is helpful whether or not
an ad is shown. Ads are served only when they are relevant, in budget, and
not already shown that session, with a last-second check right before the
reply is generated. Every sponsored mention carries the
\texttt{[Sponsored]} disclosure required by the FTC.
Sensitive conversations, around health, grief, or financial hardship are
excluded from advertising at the earliest possible point in the pipeline,
and no bid or budget can override that. And the whole thing runs within a
three-second response budget, falling back gracefully to a normal,
unsponsored reply if any part of the ad pipeline fails. Our working prototype
runs end-to-end through each of these behaviours,
from successful ad placement and sensitivity suppression through to adaptive
context revisions, image generation with product placement, budget
exhaustion, and URL-based ad drafting. To the best of our knowledge, no
existing open system brings semantic ad matching, natural LLM-written
integration with mandatory disclosure, sensitivity-aware gating,
multi-dimensional engagement measurement and organic product placement in generated images
together within a single platform. Our optimiser also operates 
without human prompting, surfacing its own diagnoses and applying its own fixes 
within guardrails the platform enforces.

\end{abstract}